% ============================================================================
%  PRELIMINARIES -- drop-in section for the ASIA CCS / ACM CCS paper.
%  Ported and trimmed from VLI_report.tex (Sections "Notation",
%  "Exact identities", "Baseline estimators"), per instruction: this section
%  exists only to support Method, not to stand alone. All proofs are moved
%  to the appendix; only statements and formulas Method cites are kept.
%
%  Dependencies this file assumes the master document's preamble provides
%  (all already present in VLI_report.tex's preamble):
%    \newtheorem{lemma}{Lemma}, {theorem}{Theorem}, {proposition}{Proposition}
%    \newcommand{\pos}[1]{\left[#1\right]_+}, \simplex, \softmax, \LSE,
%    \argmin, \argmax, \R, \E, \one
%    \usepackage{array}  -- for well-behaved p{} column vertical alignment
%    in the notation table below, sized to \columnwidth so it fits one
%    column under a two-column (ACM) class and the full text width under
%    a one-column one.
%    The baseline-comparison table uses \begin{table*}, which spans both
%    columns under a two-column class (reduces to an ordinary table under
%    a one-column one) -- it is a float, so it may drift to the top of a
%    page/column rather than sit exactly where called.
%  \label{app:proofs} is a forward reference to the appendix holding the
%  proofs of Lemma~\ref{lem:norm} and Lemma~\ref{lem:identities}; bind it
%  when this section is placed in the full paper. \ref{sec:threat} and
%  \ref{sec:method} likewise point at sections defined elsewhere.
%
%  Cut relative to VLI_report.tex: both lemma proofs; the "difficulty
%  spread" row $\kappa$ and the confusion-matrix row $Q$, $v$ (kept only if
%  Method's generalized estimator / oracle bound is included -- add back
%  the two table rows and the $Q$ definition display if so); the long
%  discussion paragraphs around each lemma; the three-part zero-sum remark
%  is reduced to the one consequence Method uses.
% ============================================================================

\section{Preliminaries}
\label{sec:prelim}

\subsection{Setting and notation}
\label{sec:prelim-notation}

Recall from Section~\ref{sec:threat} that the adversary observes one
client's accumulated final-layer gradient over a window of $m$ local
steps of batch size $B$. We analyze this window as $n = mB$ examples drawn
from $C$ classes. For example $i$:

\begin{center}
\small
\renewcommand{\arraystretch}{1.15}
\begin{tabular}{@{}l p{0.72\columnwidth}@{}}
\toprule
Symbol & Meaning \\
\midrule
$y_i$ & true label of example $i$, $y_i \in \{1,\dots,C\}$ \\
$z_i$ & logits, $z_i = W e_i + b$, with embedding $e_i$ \\
$p_i$ & predicted class distribution, $\softmax(z_i)$ \\
$\ell_i$ & cross-entropy loss, $-\log p_{i,y_i}$ \\
$L_{\mathrm{train}}$ & mean training loss on the window, $n^{-1}\sum_i \ell_i$ \\
\midrule
$N_j$ & true count of class $j$; $\sum_j N_j = n$ \\
$P_j$ & total predicted mass on class $j$, $\sum_i p_{i,j}$ \\
$G$ & \textbf{bias-gradient signal}, $\sum_{i=1}^n (p_i - e_{y_i})$ \\
\midrule
$\epsilon_i$ & per-example error, $1 - p_{i,y_i}$ \\
$\bar\epsilon$ & mean error on the window, $n^{-1}\sum_i \epsilon_i$ \\
$a_j$ & \textbf{error mass} of class $j$, $\sum_{i:y_i=j}\epsilon_i$ (zero if absent) \\
\bottomrule
\end{tabular}
\end{center}

A quantity carrying a bar or a hat ($\bar p$, $\bar O$, $\mathrm{shift}$,
$\mathrm{offset}$) is measured on an auxiliary dataset available to the
attacker in the sense of Section~\ref{sec:threat}'s side-data tier, never
on the attacked window itself.

\paragraph{Two gradient branches.}
$G$ above is the \emph{bias}-gradient signal; it is not the only quantity
these attacks read. Writing $o_i \in \R^D$ for example $i$'s embedding
(the classifier's input) and $s_i = \one^\top o_i$ for its coordinate
sum, the final layer's \emph{weight} gradient summed over the embedding
dimension is
\[
 \widetilde G_j \;=\; \sum_{d=1}^D \nabla_{W_{jd}}
 \;=\; \sum_{i=1}^n (p_{i,j} - [y_i=j])\,s_i ,
\]
an $s_i$-weighted version of the same per-example residual that defines
$G$. In the architectures we attack, the embedding is taken after a ReLU
(the final residual block's own activation for ResNet, the penultimate
ReLU for VGG), so $o_i \ge 0$ coordinatewise and hence $s_i \ge 0$: the
weighting is a genuine nonnegative average, not a sign-flipping
distortion, but it is a distortion whenever $s_i$ varies across examples.
LLG, LLG$^+$, and ZLG$^+$ read $\widetilde G$ (the implementation sums the
weight gradient over the feature dimension); RLU and \ourmethod{} read
$G$ directly from the bias gradient. Only the second branch inherits
Lemma~\ref{lem:norm}'s bound unconditionally; for the first,
$\|\widetilde G\|_1 \le (\max_i s_i)\,\|G\|_1$ in the worst case, and
Section~\ref{sec:method} returns to what this costs ZLG$^+$ in
particular, whose own calibration divides by an auxiliary estimate
$\bar O$ of $\mathrm{mean}_i\, s_i$ for exactly this reason.

\paragraph{Signal convention.}
An implementation accumulates per-minibatch \emph{mean}-loss gradients.
With all minibatches of size $B$ and $m$ steps per window, the accumulated
signal differs from the example-sum $G$ above by a common positive
factor, and possibly by a further common window normalization. Every
formula below that we mark scale-invariant is unaffected by such a
factor. This is a separate concern from Section~\ref{sec:threat}'s note on
reading a weight delta as a gradient sum: that note is about whether the
\emph{signal} is read correctly at all; this one is about a harmless
constant once it is.

\subsection{Two facts used throughout}
\label{sec:prelim-facts}

The following hold per example, and therefore for any window regardless
of how its $C$ classes are distributed; proofs are in
Appendix~\ref{app:proofs}.

\begin{lemma}[Gradient norm is controlled by training error]
\label{lem:norm}
For every example, $\|p_i - e_{y_i}\|_1 = 2\epsilon_i$. Consequently
\[
 \|G\|_1 \;\le\; 2n\bar\epsilon \;\le\; 2n L_{\mathrm{train}} .
\]
\end{lemma}

This is the quantitative statement of the phenomenon Section~\ref{sec:method}
turns into a collapse argument: as the model fits its training data, the
entire signal available to any attacker shrinks proportionally to the
training error. It does not say the label information disappears --
Lemma~\ref{lem:identities} shows it survives in the \emph{direction} of
$G$ -- only that any procedure comparing $G$ against a fixed external
magnitude will eventually be dominated by that magnitude.

\begin{lemma}[Count and error-mass identity]
\label{lem:identities}
$G_j = P_j - N_j$, and if $N_j = 0$ then $G_j = P_j \ge 0$. Moreover
\[
 \sum_{j=1}^C G_j = 0 .
\]
\end{lemma}

We use the zero-sum property once, directly: a nonzero $G$ must have a
strictly positive coordinate, so the anchor \ourmethod{} peels against in
Section~\ref{sec:method} always exists whenever the signal is not
identically zero.

\subsection{Baseline estimators, in unified notation}
\label{sec:prelim-baselines}

We write each published baseline surveyed in Section~\ref{sec:related} in
the form in which it is implemented, so that Section~\ref{sec:method} can
identify precisely which quantity each one borrows from outside the
attacked window. Throughout, $g = G/n$ denotes the normalized signal.

\paragraph{iLRG (Ma et al., 2023).}
Recovers a class-wise embedding estimate as the ratio of the weight and
bias gradients, $\hat e_j = \nabla_{W_j}/\nabla_{b_j}$, feeds it through
the classifier to obtain recovered probability vectors, and solves a
linear system for the counts by pseudo-inverse.

\paragraph{LLG (Wainakh et al., 2021).}
Reads the weight-gradient branch $\widetilde G$ (Section~\ref{sec:prelim-notation}),
and computes a single scalar \emph{impact} from it,
\[
 \mathrm{impact} = \Bigl(1 + \tfrac1C\Bigr)\frac{\sum_{j : \widetilde G_j<0} \widetilde G_j}{B},
\]
selects every class with a negative coordinate, and then greedily assigns
the remaining slots to the current $\argmin$ coordinate, decrementing by
$\mathrm{impact}$ after each assignment.

\paragraph{LLG$^+$ (Wainakh et al., 2021).}
As LLG, but first subtracts an auxiliary-measured offset from the same
branch,
\[
 \widetilde G \longleftarrow \widetilde G - \mathrm{offset}, \qquad
 \mathrm{offset}_j = \text{mean non-target coordinate $j$ on auxiliary data},
\]
and averages the auxiliary impact with the one measured before training.

\paragraph{ZLG (Geng et al., 2021; original, dummy-data form).}
Also reads $\widetilde G$, and compares it against a per-class confidence
$\bar p$ and the embedding-sum estimate $\bar O$ introduced above:
\begin{equation}
 \widehat N_j = n\Bigl(\bar p_j - \frac{\widetilde G_j}{\bar O}\Bigr),
 \label{eq:zlg}
\end{equation}
followed by clipping at zero and a multiplicative rescaling to total $n$;
dividing by $\bar O$ is this method's own attempt to undo the
embedding-norm weighting on $\widetilde G$. In its original form, both
$\bar p$ and $\bar O$ are estimated by passing \emph{dummy},
synthetically generated images through the current model -- no real
dataset is consulted, consistent with Section~\ref{sec:related}'s account
of this method as self-contained.

\paragraph{ZLG$^+$ (real-auxiliary-data form, Chen \& Vikalo, 2024).}
Identical to \eqref{eq:zlg}, except $\bar p$ and $\bar O$ are instead
estimated by passing real auxiliary samples through the current model.
As Section~\ref{sec:related} notes, this substitution is Chen \&
Vikalo's, introduced as a fairness baseline for RLU, not Geng et al.'s own
design.

\paragraph{RLU (Chen and Vikalo, 2024).}
\begin{equation}
 \widehat N = n \cdot \frac{\pos{\mathrm{shift}\cdot\bar\rho - G}}
                           {\bigl\|\pos{\mathrm{shift}\cdot\bar\rho - G}\bigr\|_1},
 \qquad \mathrm{shift} = \softmax(\mu_{\mathrm{aux}}),
 \label{eq:rlu}
\end{equation}
where $\mu_{\mathrm{aux}}$ is the vector of mean logits per class measured
on auxiliary data and $\bar\rho$ is the rate-knowledge scalar of
Section~\ref{sec:threat}; a posterior search then perturbs $\widehat N$ in
units of $m$ to better match $\mu_{\mathrm{aux}}$.

\begin{table*}[t]
\centering
\caption{Which gradient branch and which external quantity each baseline
reads, in the unified notation of this section.}
\label{tab:baselines}
\begin{tabular}{@{}lllll@{}}
\toprule
Method & Gradient branch & External quantity & Units & Form \\
\midrule
iLRG      & $\nabla W/\nabla b$ (ratio) & none            & ---                                        & pseudo-inverse \\
LLG       & $\widetilde G$              & none            & ---                                        & self-calibrated greedy \\
LLG$^+$   & $\widetilde G$              & $\mathrm{offset}$ & absolute, $\propto \bar\epsilon_{\mathrm{aux}}$ & additive \\
ZLG       & $\widetilde G$              & $\bar p$, $\bar O$ & absolute, dummy-data estimate            & \eqref{eq:zlg} \\
ZLG$^+$   & $\widetilde G$              & $\bar p$, $\bar O$ & absolute, $\propto \bar\epsilon_{\mathrm{aux}}$ & \eqref{eq:zlg} \\
RLU       & $G$                         & $\mathrm{shift}$  & absolute, $\propto \bar\epsilon_{\mathrm{aux}}$ & \eqref{eq:rlu} \\
\ourmethod{} & $G$                      & none            & ---                                        & Section~\ref{sec:method} \\
\bottomrule
\end{tabular}
\end{table*}

Section~\ref{sec:method} shows that every row depending on an absolute
external quantity collapses as $\bar\epsilon \to 0$, and that this is not
a coincidence of these designs but a consequence of Lemma~\ref{lem:norm}
applied to whatever such quantity each one chooses. The gradient-branch
column matters separately: for the four methods reading $\widetilde G$
(LLG, LLG$^+$, ZLG, ZLG$^+$), even the \emph{signal} side of the
comparison carries the extra embedding-norm factor introduced above, so
Lemma~\ref{lem:norm}'s clean bound transfers to them only once that
factor is accounted for. ZLG is listed here despite reading no real
external dataset because $\bar p$ and $\bar O$ are still quantities
measured \emph{outside the attacked window} -- on dummy images rather
than on the window's own examples -- which is the property
Theorem~\ref{thm:collapse} in Section~\ref{sec:method} actually turns on,
not whether that outside source is real data.
